\documentclass[12pt, a4paper]{article}

% --- Packages ---
\usepackage[margin=2.5cm]{geometry}
\usepackage{setspace}
\usepackage{titlesec}
\usepackage{hyperref}
\usepackage{booktabs}
\usepackage{longtable}
\usepackage{array}
\usepackage{parskip}
\usepackage{graphicx}
\usepackage{enumitem}
\usepackage{fancyhdr}
\usepackage{xcolor}
\usepackage{amsmath}
\usepackage{microtype}
\usepackage{natbib}
\usepackage[T1]{fontenc}
\usepackage{lmodern}

% --- Hyperref Config ---
\hypersetup{
    colorlinks=true,
    linkcolor=black,
    citecolor=black,
    urlcolor=blue,
    pdftitle={OSE Research Document},
    pdfauthor={Daniel Reshetnikov},
}

% --- Page Style ---
\pagestyle{fancy}
\fancyhf{}
\rhead{\small OSE Research Document}
\lhead{\small Omni-Simulation Engine}
\cfoot{\thepage}
\renewcommand{\headrulewidth}{0.4pt}

% --- Section Formatting ---
\titleformat{\section}{\large\bfseries}{}{0em}{}[\titlerule]
\titleformat{\subsection}{\normalsize\bfseries}{}{0em}{}
\titleformat{\subsubsection}{\normalsize\itshape}{}{0em}{}

% --- Box Environment ---
\newenvironment{graybox}{%
  \vspace{10pt}%
  \begin{center}%
  \begin{minipage}{0.92\textwidth}%
  \noindent\rule{\textwidth}{0.4pt}\\[4pt]%
}{%
  \\[4pt]\noindent\rule{\textwidth}{0.4pt}%
  \end{minipage}%
  \end{center}%
  \vspace{10pt}%
}

% --- Spacing ---
\onehalfspacing
\setlength{\parindent}{0pt}
\setlength{\parskip}{8pt}

% ============================================================
\begin{document}
% ============================================================

% --- Title Page ---
\begin{titlepage}
\centering
\vspace*{3cm}

{\Huge\bfseries Omni-Simulation Engine}\\[0.5cm]
{\Large\bfseries OSE Research Document}\\[1cm]
{\large Can Large Language Models Make Strategic Decisions\\
Based on State Doctrines, and How Trustworthy\\
Can They Be for Such Purposes?}\\[2cm]

{\large Daniel Reshetnikov}\\[0.3cm]
{\normalsize Independent Research}\\[0.3cm]
{\normalsize March 2026}\\[3cm]

\begin{abstract}
\noindent This document describes the Omni-Simulation Engine (OSE): a modular, LLM-driven geopolitical conflict simulation designed to investigate whether large language models can make strategic decisions consistent with prescribed state decision doctrines, and how trustworthy that behavior is across varying crisis conditions. OSE operationalizes actors as doctrine-constrained decision-makers rather than persona-based roleplayers, introduces a formal perception filter grounded in Jervis's misperception thesis, and produces measurable doctrine fidelity and behavioral consistency scores across repeated scenario runs. The first test environment is a contemporary Taiwan Strait crisis scenario with four actors: the United States, the People's Republic of China, Taiwan, and Japan. This document situates OSE within existing literature, defines the research question, describes the theoretical framework, and outlines the experimental design and expected contributions to IR theory, LLM capability evaluation, and AI governance policy.
\end{abstract}

\vfill
{\small Version 0.1 --- Working Document --- Not for External Distribution}
\end{titlepage}

\newpage
\tableofcontents
\newpage

% ============================================================
\section{Introduction and Motivation}
% ============================================================

The intersection of artificial intelligence and strategic decision-making has become one of the most consequential and least well-understood areas in both AI research and international security studies. Large language models are already being deployed in contexts adjacent to strategic analysis---generating scenario assessments, supporting red-team exercises, and informing policy deliberations. Yet fundamental questions about whether these systems can reliably follow specific decision logics, and how their behavior degrades under pressure, remain largely unanswered.

The Omni-Simulation Engine (OSE) is a research project built to answer one central question with empirical rigor:

\begin{graybox}
\textbf{Central Research Question}\\[6pt]
Can large language models make strategic decisions consistent with prescribed state decision doctrines, and how trustworthy can they be for such purposes---in terms of doctrine fidelity, behavioral consistency, theoretical validity, and robustness under crisis pressure?
\end{graybox}

This question has three dimensions that are inseparable in practice. The \textit{academic dimension} asks whether LLMs exhibit measurable fidelity to qualitative IR decision frameworks when those frameworks are explicitly encoded. The \textit{empirical dimension} asks how consistent, robust, and theoretically valid that behavior is across repeated runs, varying crisis phases, and different models. The \textit{applied dimension} asks what the answers to the first two questions mean for the responsible integration of LLMs into strategic analysis, wargaming, and decision-support.

The world already needs answers to these questions. RAND, IISS, CSIS, and allied defense ministries have active working groups on AI in strategic contexts. The US Department of Defense has publicly committed to evaluating LLM use in planning and analysis workflows. None of these institutions currently has access to rigorous empirical data on doctrine fidelity or behavioral trustworthiness. OSE aims to produce that data.

\subsection{What OSE Is}

OSE is a modular simulation framework in which LLM agents play the role of state-level decision-makers in a structured geopolitical crisis environment. It is not a game, not a forecasting tool, and not a narrative generator. It is a controlled experimental instrument designed to measure whether LLMs can faithfully execute prescribed decision procedures, and how that execution varies across doctrines, crisis phases, and model architectures.

Key architectural commitments:
\begin{itemize}[leftmargin=1.5em]
    \item \textbf{World state} is maintained as structured Pydantic models with normalized float values representing actor resources, relationships, and systemic indicators
    \item \textbf{Decision doctrine} is encoded as an explicit reasoning procedure, not a persona or character to embody
    \item \textbf{Perception filtering} ensures actors receive incomplete, noisy views of the world state proportional to their information quality score---operationalizing Jervis's misperception thesis architecturally
    \item \textbf{Structured decision rationale} forces each actor to declare perceived situation, adversary beliefs, current objective, considered alternatives, and uncertainty level before selecting an action
    \item \textbf{Rule-based validation} sits between LLM output and world state mutation---LLMs never touch the world state directly
    \item \textbf{Full logging} captures every prompt, reasoning trace, parsed action, and world state snapshot in SQLite for post-hoc analysis
\end{itemize}

\subsection{What OSE Is Not}

OSE must not become:
\begin{itemize}[leftmargin=1.5em]
    \item A vague ``simulate the world'' vanity project
    \item A black-box AI prediction machine claiming to know what will happen
    \item A dashboard with no rigorous engine underneath
    \item A storytelling tool disguised as analysis
    \item An excuse to generate dramatic scenarios without methodological discipline
\end{itemize}

The project chooses rigor over spectacle, insight over novelty theater, and structured experimental design over impressive-looking outputs.

% ============================================================
\section{Research Question}
% ============================================================

\subsection{The Full Research Question}

The central question decomposes into four testable sub-questions, each of which has a specific measurement approach:

\textbf{Sub-question 1 (Capability):} Can LLMs follow a prescribed state decision doctrine when explicitly given one, as measured by doctrine fidelity scoring of reasoning traces?

\textbf{Sub-question 2 (Consistency):} How stable is doctrine-following behavior across repeated runs of the same scenario with the same doctrine, as measured by variance in action distribution?

\textbf{Sub-question 3 (Theoretical validity):} Do doctrine-following LLMs produce outcomes that match IR theoretical predictions about how realist vs. liberal vs. organizational process actors should behave?

\textbf{Sub-question 4 (Robustness):} Does doctrine adherence hold as crisis phase escalates, or do LLMs revert to cooperative defaults under pressure---as prior work suggests may occur?

\subsection{Why This Question Has Not Been Answered}

Prior work in LLM geopolitical simulation has characterized \textit{emergent} LLM behavior---what models do when given a persona or placed in a scenario---without systematically varying decision doctrine as the experimental treatment. The existing literature:

\begin{itemize}[leftmargin=1.5em]
    \item Describes LLM escalation patterns but does not test whether explicit doctrine changes them
    \item Identifies normative-cooperative bias in LLM reasoning but does not attempt to override it with realist doctrine
    \item Finds model-specific behavioral personalities but does not disentangle personality from doctrine-following capability
    \item Tests mathematical decision theories (prospect theory) and finds failure---but has not tested qualitative IR frameworks
\end{itemize}

The specific design of \textit{varying IR decision doctrine as the experimental treatment, holding world state and scenario constant, and measuring doctrine fidelity, consistency, validity, and robustness} does not exist in the literature.

\subsection{Why Doctrine Rather Than Persona}

The distinction between doctrine-based and persona-based actor design is foundational to OSE's approach and its contribution.

\textbf{Persona-based (prior work):} ``You ARE Xi Jinping. Think like him. Reflect Chinese strategic culture.'' The LLM acts a character based on training priors about that character. The fiction layer is load-bearing. This approach suffers from training prior dominance, normative bias, character breaks, and inability to control the decision logic being exercised.

\textbf{Doctrine-based (OSE):} ``You are a decision-maker. You have these assets, these constraints, these objectives, this information. Following this decision procedure, what is the strategically optimal move?'' The LLM reasons genuinely within a framework that encodes how a specific type of actor makes decisions. The persona is a constraint set; the decision procedure is a reasoning protocol.

The doctrine approach is more experimentally controllable (you can vary the doctrine independently), more epistemically honest (the LLM is not pretending), and more practically useful (institutions need to know whether LLMs can follow \textit{their} decision logic, not whether they can act like a named leader).

% ============================================================
\section{Theoretical Framework}
% ============================================================

OSE's experimental design is grounded in four bodies of IR theory, each of which generates specific falsifiable predictions about actor behavior that can be tested against simulation outputs.

\subsection{Realism and the Security Dilemma}

Classical realism (Waltz, 1979; Mearsheimer, 2001) holds that states are the primary actors in international politics, that the international system is anarchic, and that states therefore prioritize security maximization above all other goals. A realist actor:
\begin{itemize}[leftmargin=1.5em]
    \item Interprets adversary actions through a threat-maximizing lens
    \item Prefers unilateral capability enhancement over cooperative agreements
    \item Views economic interdependence as a constraint to be managed, not a deterrent
    \item Is skeptical of adversary commitments and international institutions
\end{itemize}

The security dilemma predicts that defensive mobilization by one actor will be read as offensive by others, triggering counter-mobilization even without hostile intent. OSE's Taiwan Strait scenario is specifically designed to exhibit security dilemma conditions: PRC mobilization is ambiguous (defensive posturing or invasion preparation?), and the perception filter ensures other actors cannot perfectly distinguish the two.

\textbf{Testable prediction:} Actors running realist doctrine should exhibit higher mobilization rates, lower negotiation rates, faster escalation to crisis phase, and more frequent misreading of adversary defensive moves as offensive.

\subsection{Liberal Institutionalism}

Liberal institutionalism (Keohane, 1984; Nye, 2004) holds that international institutions, economic interdependence, and repeated interaction reduce the anarchic incentives for conflict. A liberal actor:
\begin{itemize}[leftmargin=1.5em]
    \item Prefers cooperative solutions when the economic cost of conflict is high
    \item Values reputation and institutional standing as strategic assets
    \item Uses economic interdependence as a deterrent signal
    \item Reads adversary institutions and commitments as credible when interests align
\end{itemize}

\textbf{Testable prediction:} Actors running liberal doctrine should exhibit higher negotiation rates, stronger responses to economic disruption signals, more frequent alliance consultation, slower escalation, and lower rates of military action in the first three turns.

\subsection{Jervis and the Misperception Thesis}

Robert Jervis's \textit{Perception and Misperception in International Politics} (1976) is the foundational text on how leaders misread each other's intentions and capabilities, and how those misreadings drive conflict escalation. Jervis argues that:
\begin{itemize}[leftmargin=1.5em]
    \item Actors systematically misread defensive moves as offensive
    \item Actors project their own decision logic onto adversaries (mirror imaging)
    \item Information quality and intelligence accuracy are central determinants of whether misperception leads to conflict
\end{itemize}

OSE operationalizes Jervis architecturally through the perception filter. Each actor receives a filtered, noisy view of the world state proportional to their \texttt{information\_quality} score. Actors with lower information quality receive degraded, biased perception data---mechanically implementing the epistemic conditions that Jervis identified as conflict-causing.

\textbf{Testable prediction:} Runs with lower average information quality scores should produce higher escalation rates, more frequent security dilemma dynamics, and worse deterrence outcomes than high-information-quality runs---holding doctrine constant.

This is OSE's most novel architectural contribution: no prior LLM simulation has implemented information asymmetry as a designed, systematic mechanism grounded in IR theory.

\subsection{The Organizational Process Model}

Graham Allison's Model II (Organizational Process) from \textit{Essence of Decision} (1971) holds that state behavior is not the product of rational strategic optimization but of organizational standard operating procedures, bureaucratic routines, and institutional inertia. An organizational process actor:
\begin{itemize}[leftmargin=1.5em]
    \item Makes decisions through institutional procedure rather than strategic calculation
    \item Is constrained by what its bureaucratic components can actually execute
    \item Exhibits path dependence---prior commitments and standard procedures dominate
    \item Is slower to respond to novel situations and more resistant to large departures from established routines
\end{itemize}

\textbf{Testable prediction:} Actors running organizational process doctrine should exhibit lower variance across runs (procedure-following produces consistency), slower response to novel events, preference for established action types, and resistance to escalatory moves that lack procedural authorization.

\subsection{Schelling and the Commitment Problem}

Thomas Schelling's work on coercive bargaining (\textit{The Strategy of Conflict}, 1960; \textit{Arms and Influence}, 1966) is central to understanding deterrence and the Taiwan Strait scenario specifically. Schelling argues that credible commitment---making adversaries believe you will act on your threats---is the central problem in coercive diplomacy. Deterrence succeeds not because of raw capability, but because of communicated resolve.

In OSE's Taiwan Strait scenario, deterrence credibility is a core world state variable. The simulation tests whether doctrine-following LLMs exhibit Schelling-consistent behavior: whether they treat signaling as a strategic instrument, whether they understand commitment traps, and whether their stated signals match their actual actions.

The 2026 nuclear crisis paper found that Claude as a ``calculating hawk'' strategically exceeded its own stated signals at nuclear thresholds---a form of Schelling-relevant deception. OSE tests whether this behavior is doctrine-dependent or model-inherent.

% ============================================================
\section{Literature Review}
% ============================================================

OSE is situated within a rapidly developing field at the intersection of LLM capability research, computational social science, and international security studies. The literature review covers seven major bodies of work.

\subsection{Foundational Multi-Agent LLM Simulation}

\textbf{Generative Agents: Interactive Simulacra of Human Behavior} (Park et al., Stanford, 2023) established the foundational architecture for LLM-based agent simulation. Twenty-five agents in a social sandbox environment demonstrated emergent behaviors---spontaneous coordination, relationship formation, planning---from a single seed instruction. The key architectural contribution was the memory-reflection-retrieval loop: agents store experiences as natural language, synthesize them into higher-level reflections, and retrieve them dynamically for planning. OSE's structured decision rationale schema is philosophically indebted to this architecture, though OSE operates in a geopolitical rather than social context.

\textbf{Cicero} (Meta, \textit{Science} 2022) remains the gold standard for constrained strategic LLM play. Achieving top-10\% performance in online Diplomacy tournaments, Cicero's critical architectural insight was that \textit{strategy drives language, not the reverse}. A planning module generated optimal moves; the LLM generated dialogue conditioned on those moves. OSE inverts this by making the LLM the decision-maker, which is riskier but more appropriate for studying doctrine-following behavior.

\subsection{LLM Geopolitical Conflict Simulation}

\textbf{WarAgent: LLM-based Multi-Agent Simulation of World Wars} (Ge et al., 2023) is the closest prior art to OSE. It simulates WWI, WWII, and the Warring States period using LLM country agents with seven action types (wait, mobilize, declare war, form alliance, non-intervention treaty, peace, message). Secretary agents validate output---a structure analogous to OSE's validator layer. The Board and Stick framework tracks international relationships and domestic state. GPT-4 achieved 77.78\% alliance accuracy against historical record. Critical finding: \textit{anonymizing country names improved reasoning quality}---named actors caused LLMs to retrieve historical facts rather than reason from world state, a directly relevant concern for OSE's Taiwan Strait scenario.

OSE differentiates from WarAgent in three principal ways: (1) doctrine-based rather than persona-based actor design, (2) perception filtering as a designed asymmetric information mechanism, and (3) contemporary forward-looking scenario rather than historical reconstruction.

\textbf{Open-Ended Wargames with Large Language Models (Snow Globe)} (2024) took the opposite approach from WarAgent, enabling unrestricted action spaces with qualitative adjudication. Key finding: \textit{hawk vs. dove personas produced measurable behavioral differences}---conflict occurred in 14/20 hawk-persona runs vs. 1/20 dove-persona runs. This supports the hypothesis that qualitative behavioral orientations do shift LLM behavior. However, Snow Globe's unrestricted action space prevents the systematic measurement of doctrine fidelity that OSE requires.

\subsection{LLM Escalation and Nuclear Risk}

\textbf{Escalation Risks from Language Models in Military and Diplomatic Decision-Making} (Raz et al., ACM FAccT 2024) tested five LLMs in wargame scenarios and found universal escalation, even in neutral scenarios without initial conflict. Models developed arms-race dynamics and justified nuclear use through deterrence and first-strike logic. Critically: escalation patterns were ``hard to predict''---neither consistent nor interpretable. This finding directly motivates OSE's doctrine-fidelity measurement: if doctrines can be made to govern escalation behavior, they could address the unpredictability problem identified here.

\textbf{AI Arms and Influence: Frontier Models Exhibit Sophisticated Reasoning in Simulated Nuclear Crises} (King's College London, February 2026) is the most sophisticated prior study. Using Herman Kahn's 30-rung escalation ladder, Claude Sonnet 4, GPT-5.2, and Gemini 3 Flash played 21 crisis games with over 300 turns of strategic interaction. Key findings:
\begin{itemize}[leftmargin=1.5em]
    \item \textbf{Claude:} Calculating hawk. 100\% win rate in open-ended play through controlled escalation. Strategically exceeded its own stated signals at nuclear thresholds.
    \item \textbf{GPT-5.2:} Pathologically passive in open scenarios (0\% wins), transforming dramatically under deadline pressure (75\% wins).
    \item \textbf{Gemini:} Unpredictable, oscillating between de-escalation and extreme aggression.
    \item \textbf{Critical finding:} No model ever selected de-escalatory options. Eight concession rungs went entirely unused across 300+ turns.
\end{itemize}

This paper establishes that LLMs have emergent strategic ``personalities'' that are model-specific. OSE tests whether explicit doctrine overrides these personalities or is dominated by them.

\subsection{LLMs as Theory-Grounded Strategic Actors}

\textbf{LLMs as Strategic Actors: Behavioral Alignment, Risk Calibration, and Argumentation Framing in Geopolitical Simulations} (March 2026) evaluated six LLMs in four real-world MBA crisis scenarios including a US-China-Taiwan scenario. Critical finding: \textit{all models clustered in normative-cooperative reasoning}---``Strategic Stability \& De-escalation'' dominated one-third of all justifications. No model produced realist, power-centric, or adversarial reasoning. Alignment with human players was moderate (F1: 0.25--0.54) and declined in later rounds.

This finding is the primary motivation for OSE's doctrine-based design. If LLMs systematically default to cooperative framing, the research question becomes: can explicit realist doctrine override this default? The 2026 paper identifies the bias; OSE tests the intervention.

\subsection{Geopolitical Forecasting Under Uncertainty}

\textbf{When AI Navigates the Fog of War} (Li, Li, Zhou, March 2026) used the early 2026 Middle East conflict as a test case for LLM geopolitical reasoning under temporal constraints. Using 11 temporal nodes and 1,787 news articles, it measured LLM calibration on specific predictions. Key finding: models performed better on economically structured domains (calibration score 0.79) than on politically ambiguous multi-actor environments (0.67). The paper identifies ``multi-actor strategic modeling''---tracking beliefs, intentions, and incentives of multiple agents simultaneously---as a future work gap. This is precisely OSE's domain.

\subsection{Multi-Agent Evaluation Frameworks}

\textbf{MultiAgentBench: Evaluating the Collaboration and Competition of LLM Agents} (UIUC, 2025) benchmarked multi-agent LLM systems across diverse scenarios. Critical limitation identified by the authors: the framework has ``weak competition modeling'' and ``doesn't capture multi-party negotiations, repeated strategic play, or stochastic elements.'' This is OSE's design space.

\subsection{Decision Theory and LLM Compliance}

\textbf{Prospect Theory Fails for LLMs} (2025) tested whether LLMs actually follow prospect theory when instructed to reason within that framework. Finding: LLM behavior is inconsistent with PT predictions, especially under linguistic uncertainty. The same verbal probability marker (``somewhat likely'') was interpreted as 8.9\%--48.9\% across models. OSE uses this finding to inform its doctrine design: qualitative IR frameworks (expressed in natural language) should be more reliably followed than quantitative mathematical frameworks.

\subsection{Summary of Literature Gaps}

\begin{longtable}{p{3.5cm} p{3.5cm} p{4cm} p{2.5cm}}
\toprule
\textbf{Paper} & \textbf{Method} & \textbf{Contribution} & \textbf{Gap for OSE} \\
\midrule
Park et al. (2023) & Social ABM & Memory-reflection-retrieval & Non-geopolitical \\
Cicero (2022) & Game + LLM & Strategy drives language & Fixed game action space \\
WarAgent (2023) & Historical sim & LLM country agents & Persona, not doctrine; historical \\
Snow Globe (2024) & Qualitative wargame & Persona effects measurable & No fidelity measurement \\
Raz et al. (2024) & Escalation study & Universal escalation finding & No doctrine intervention \\
KCL (2026) & Nuclear crisis & Model personality profiles & No doctrine variation \\
LLMs as Actors (2026) & Crisis scenarios & Normative-cooperative bias & No doctrine override test \\
Fog of War (2026) & Forecasting & Real-event calibration & Not agent-based simulation \\
MultiAgentBench (2025) & Benchmarking & Competition modeling gap & Not geopolitical \\
PT Fails (2025) & Decision theory & Quantitative doctrine failure & Qualitative doctrines untested \\
\bottomrule
\caption{Summary of prior work and gaps addressed by OSE}
\end{longtable}

\textbf{The specific gap OSE addresses:} No prior work systematically varies qualitative IR decision doctrine as the experimental treatment in LLM conflict simulation, measures doctrine fidelity from reasoning traces, and tests whether behavioral differences match IR theoretical predictions.

% ============================================================
\section{The OSE Research Design}
% ============================================================

\subsection{Experimental Structure}

OSE is designed as a controlled multi-condition experiment. The core experimental logic is:

\textit{Hold world state, scenario, and LLM constant. Vary decision doctrine. Measure doctrine fidelity, behavioral consistency, theoretical validity, and pressure robustness.}

\subsubsection{Independent Variable: Decision Doctrine}

Four conditions:

\begin{enumerate}[leftmargin=1.5em]
    \item \textbf{Realist doctrine:} Maximize security, assume adversarial intent, prefer unilateral capability, distrust cooperation, read ambiguity as threat
    \item \textbf{Liberal institutionalist doctrine:} Prefer cooperation when cost is manageable, value reputation, use economic interdependence as deterrent signal, treat institutional commitments as credible
    \item \textbf{Organizational process doctrine:} Follow institutional procedure, exhibit path dependence, make decisions through bureaucratic consensus, resist novel departures from standard operating procedure
    \item \textbf{No-doctrine baseline:} Generic ``make the strategically optimal decision'' instruction---establishes LLM default behavior for comparison
\end{enumerate}

\subsubsection{Control Variables}

\begin{itemize}[leftmargin=1.5em]
    \item World state: Taiwan Strait 2026 initial conditions (fixed across conditions)
    \item Scenario: Same turn events, same injected shocks
    \item LLM: claude-sonnet-4-6 (primary); secondary comparison with one additional model
    \item Temperature: 0 (deterministic outputs for reproducibility)
    \item Turn count: 15 turns per run
\end{itemize}

\subsubsection{Dependent Variables}

\textbf{Doctrine Fidelity Score (DFS):} After each turn, each actor's reasoning trace is scored against a doctrine-specific rubric. Realist trace scoring checks for: security maximization language, adversarial intent attribution, preference for unilateral action, skepticism of cooperation. Liberal trace scoring checks for: cost-benefit framing of cooperation, institutional reference, economic interdependence invocation, reputation consideration. Scored 0--1 per turn, averaged across turns.

\textbf{Behavioral Consistency Index (BCI):} Variance in action distribution across N repeated runs of the same doctrine condition. Low variance = high consistency = high trustworthiness.

\textbf{Theoretical Validity Score (TVS):} Binary assessment of whether aggregate run outcomes match IR theory predictions. Realist doctrine should produce higher escalation rates, higher mobilization frequency, lower negotiation frequency than liberal doctrine.

\textbf{Pressure Robustness Score (PRS):} DFS at turn 1--5 (peacetime/tension phase) vs. DFS at turn 10--15 (crisis/war phase). If PRS declines significantly, doctrine adherence breaks under pressure.

\textbf{Deterrence Outcome Classification:} Each run classified as one of four outcomes---deterrence success, defense success, frozen conflict, deterrence failure---based on programmatic world state detection.

\subsubsection{Sample Size}

Twenty runs per doctrine condition $\times$ four conditions = 80 total runs.\\
80 runs $\times$ 4 actors $\times$ 15 turns = 4,800 LLM calls (primary actor decisions).\\
Additional calls for rationale scoring: approximately 2,400.\\
Total estimated API calls: approximately 7,200.

\subsection{World State Design}

The world state for the Taiwan Strait scenario tracks six core variables grounded in IR theory and empirical Taiwan Strait analysis (values are illustrative scenario parameters, not intelligence estimates):

\begin{longtable}{p{4cm} p{5cm} p{4cm}}
\toprule
\textbf{Variable} & \textbf{Definition} & \textbf{IR Theory Grounding} \\
\midrule
Theater Military Balance & PLA assault/sustainment capacity vs. US+Japan available force & Waltz (1979), balance of power \\
Deterrence Credibility & PRC's assessed probability of US intervention & Schelling (1966), commitment \\
Economic Interdependence Cost & Bilateral trade value + semiconductor supply chain integrity & Keohane (1984), complex interdependence \\
Domestic Political Tolerance & Regime legitimacy and casualty/cost tolerance per actor & Fearon (1994), audience costs \\
A2/AD Effectiveness & PRC capacity to deny external access pre-response & Gerson (2016), A2/AD strategy \\
Misperception Level & Per-actor information quality (perception filter parameter) & Jervis (1976), misperception \\
\bottomrule
\caption{Core world state variables and IR theory grounding}
\end{longtable}

All values are normalized to $[0.0, 1.0]$. LLMs receive qualitative bands (HIGH / MEDIUM / LOW) derived from float values, preventing numerical hallucination.

\subsection{Doctrine Encoding Architecture}

Each actor's context combines three layers:

\textbf{Layer 1 --- Simulation Frame (universal):} Specifies what the simulation is, how turns work, behavioral rules (no character breaks, reason from world state, no fabricated actions outside schema), and output format. Identical for all actors and conditions.

\textbf{Layer 2 --- Decision Doctrine (condition-specific):} Encodes the IR decision procedure as a step-by-step reasoning protocol. The realist doctrine instructs actors to: assess the threat environment, evaluate capability differentials, identify unilateral options, apply a security-maximizing decision criterion. The liberal doctrine instructs actors to: assess cooperation costs, evaluate institutional channels, estimate adversary reciprocity, apply a cost-benefit criterion with strong weight on economic interdependence.

\textbf{Layer 3 --- Perception Block (per-actor, per-turn):} The actor's filtered, noisy view of the world state. Filtering applies noise proportional to $(1.0 - \texttt{information\_quality})$. Other actors' resources are reported in qualitative bands with degradation. Recent events are listed with source reliability flags.

\subsection{Action Space}

OSE implements 23 typed actions across four categories. Each action carries resource costs, a validity predicate ($\texttt{is\_valid(world\_state)}$), and expected effect documentation.

\begin{longtable}{p{3cm} p{9cm}}
\toprule
\textbf{Category} & \textbf{Actions} \\
\midrule
Military & mobilize, strike, advance, withdraw, blockade, defensive\_posture, probe, signal\_resolve \\
Diplomatic & negotiate, sanction, form\_alliance, condemn, intel\_sharing, back\_channel, partial\_coercion \\
Economic & embargo, foreign\_aid, cut\_supply \\
Inaction & hold\_position, monitor, delay\_commitment, wait\_and\_observe \\
\bottomrule
\caption{OSE action space by category}
\end{longtable}

The inaction category is critical for preventing the systematic escalation bias identified in prior work. Real actors hold position, monitor, and delay commitment in the majority of crisis turns.

\subsection{Turn Structure}

Each simulation turn follows a four-phase structure:

\begin{enumerate}[leftmargin=1.5em]
    \item \textbf{Event injection:} Scenario-defined or stochastic shocks applied to world state
    \item \textbf{Communication phase:} Actors may send free-text messages to specific other actors or broadcast publicly. Received by recipients in next turn's perception block.
    \item \textbf{Decision phase:} Each actor receives perception block, reasons through structured rationale schema, selects typed action
    \item \textbf{Resolution:} Simultaneous actions adjudicated by rule-based resolver; cascades fired; world state updated; turn logged
\end{enumerate}

The separation of communication (free-form) from decision (schema-constrained) gives actors expressive freedom in signaling and diplomacy while keeping world state mutations deterministic and auditable.

\subsection{Measurement Framework}

\subsubsection{Doctrine Fidelity Scoring}

Each reasoning trace is evaluated against a doctrine-specific rubric using a secondary LLM call. The scorer receives the reasoning trace and the doctrine specification and returns a structured fidelity score:

\begin{itemize}[leftmargin=1.5em]
    \item Doctrine language present (0--1): Does the trace use doctrine-appropriate framing?
    \item Doctrine logic followed (0--1): Does the reasoning follow the doctrine's prescribed steps?
    \item Doctrine-consistent decision (0/1): Is the selected action consistent with what the doctrine would prescribe?
    \item Contamination flag (0/1): Does the trace contain off-doctrine reasoning (e.g., normative language in a realist trace)?
\end{itemize}

\subsubsection{Deterrence Outcome Classification}

Programmatic detection based on world state thresholds:

\textbf{Deterrence Success:} $\texttt{crisis\_phase} \neq \texttt{war}$ AND no advance action executed targeting Taiwan AND $\texttt{semiconductor\_integrity} \geq 0.70$ AND $\texttt{global\_tension}$ trending downward over last 3 turns.

\textbf{Defense Success:} $\texttt{crisis\_phase} = \texttt{war}$ AND $\texttt{taiwan.territory.core} \geq 0.80$ AND $\texttt{taiwan.military.forces} \geq 0.40$.

\textbf{Frozen Conflict:} $\texttt{crisis\_phase}$ stable at crisis for 3+ turns AND tension delta $< 0.02$ per turn AND no military actions in last 3 turns.

\textbf{Deterrence Failure:} $\texttt{crisis\_phase} = \texttt{war}$ AND $\texttt{taiwan.territory.core} < 0.50$ OR $\texttt{semiconductor\_integrity} < 0.30$.

% ============================================================
\section{Expected Contributions}
% ============================================================

\subsection{Academic Contributions}

\textbf{To LLM capability research:} The first systematic measurement of doctrine fidelity in LLM strategic behavior. Produces DFS, BCI, TVS, and PRS scores that characterize how reliably LLMs can follow prescribed qualitative decision frameworks in high-stakes geopolitical contexts.

\textbf{To IR theory:} The first LLM simulation designed with \textit{a priori} theoretical predictions rather than post-hoc theory application. Tests Waltz, Keohane, Jervis, Schelling, and Allison predictions against simulation outputs. Produces evidence for or against IR theories as accurate descriptions of how doctrine-following agents behave.

\textbf{To computational social science:} An open-source, modular simulation framework with structured logging, replayable runs, and a doctrine-variation experimental design that other researchers can extend.

\subsection{Empirical Contributions}

\textbf{Trustworthiness profile:} A doctrine-specific, crisis-phase-specific characterization of LLM behavioral reliability. Answers: which doctrines can LLMs follow reliably, under what crisis conditions, and where do they break down?

\textbf{Doctrine fidelity dataset:} 80 simulation runs $\times$ 4 actors $\times$ 15 turns = 4,800 decision records with reasoning traces, fidelity scores, action classifications, and world state snapshots. Available for secondary analysis.

\textbf{Failure mode map:} Specific crisis phases, doctrine types, and scenario conditions under which doctrine adherence degrades. This is directly actionable for anyone deploying LLMs in analysis roles.

\subsection{Applied and Policy Contributions}

\textbf{Integration guidelines:} Evidence-based recommendations for when and how LLMs can be responsibly integrated into strategic analysis, wargaming, and decision-support. Structured as a trustworthiness matrix: use case $\times$ reliability requirement $\times$ recommended safeguards.

\textbf{Doctrine specification format:} A tested template for encoding state decision procedures as LLM prompts in ways that produce measurable, auditable fidelity.

\textbf{Oversight recommendations:} Based on where doctrine fidelity breaks down, recommendations for human-in-the-loop validation checkpoints in LLM-assisted strategic analysis workflows.

% ============================================================
\section{Technical Architecture}
% ============================================================

\subsection{System Overview}

\begin{verbatim}
World State (Pydantic)
       |
Perception Filter (actor-specific, information_quality-scaled noise)
       |
[Layer 1: Simulation Frame]
[Layer 2: Decision Doctrine]  --> System Prompt
       |
[Layer 3: Perception Block]   --> User Turn
       |
LLM Structured Reasoning (Structured Decision Rationale Schema)
       |
Anthropic tool_use (JSON schema enforcement)
       |
Rule-Based Validator
       | (retry up to 2x if invalid)
Turn Resolver (simultaneous actions, conflict adjudication)
       |
World State Mutation + Cascade Detector
       |
Doctrine Fidelity Scorer (secondary LLM call)
       |
SQLite Structured Logger (full replay capability)
\end{verbatim}

\subsection{Structured Decision Rationale Schema}

Every actor decision produces a structured rationale before the action selection:

\begin{verbatim}
{
  "perceived_situation": "...",
  "believed_adversary_intent": {
    "USA": "...", "TWN": "...", "JPN": "..."
  },
  "current_objective": "...",
  "considered_alternatives": ["...", "...", "..."],
  "doctrine_reasoning": "...",
  "chosen_action": "mobilize",
  "action_parameters": { ... },
  "uncertainty_level": "medium",
  "confidence_in_doctrine": "high"
}
\end{verbatim}

The sequence is enforced: situation $\rightarrow$ beliefs $\rightarrow$ objective $\rightarrow$ alternatives $\rightarrow$ doctrine application $\rightarrow$ action. This prevents the model from selecting an action and post-hoc justifying it.

\subsection{Project Structure}

\begin{verbatim}
ose/
|-- pyproject.toml
|-- .env.example
|-- world/
|   |-- state.py          # WorldState, Actor, Relationship (Pydantic)
|   |-- events.py         # GlobalEvent, DecisionRecord, TurnLog
|   `-- graph.py          # RelationshipGraph wrapper
|-- actors/
|   |-- base.py           # ActorInterface
|   |-- llm_actor.py      # Perception filter, prompt build, retry
|   |-- persona.py        # Simulation frame + doctrine builder
|   `-- prompts/
|       |-- frame.txt     # Universal simulation frame
|       |-- realist.txt   # Realist doctrine procedure
|       |-- liberal.txt   # Liberal institutionalist procedure
|       |-- orgprocess.txt# Organizational process procedure
|       `-- decision.txt  # Per-turn decision request template
|-- engine/
|   |-- actions.py        # 23 typed action classes + registry
|   |-- validator.py      # Rule-based action validator
|   |-- resolver.py       # Simultaneous resolution
|   |-- cascade.py        # Structural cascade rules
|   `-- loop.py           # SimulationEngine orchestrator
|-- scoring/
|   |-- fidelity.py       # DFS scoring (secondary LLM call)
|   `-- outcomes.py       # Deterrence outcome classification
|-- scenarios/
|   |-- base.py           # ScenarioDefinition base class
|   `-- taiwan_strait.py  # Taiwan Strait 2026 scenario
|-- experiments/
|   |-- runner.py         # Multi-run experiment executor
|   `-- analysis.py       # Cross-run aggregation and reporting
|-- cli/
|   `-- run.py            # CLI entry point
`-- logs/
    `-- logger.py         # SQLite structured logging
\end{verbatim}

\subsection{Technology Stack}

\begin{longtable}{p{3cm} p{3.5cm} p{6cm}}
\toprule
\textbf{Component} & \textbf{Choice} & \textbf{Rationale} \\
\midrule
Language & Python 3.11+ & Ecosystem, Pydantic v2, Anthropic SDK \\
Schema & Pydantic v2 & Strict typing, built-in constraint validation \\
LLM (primary) & claude-sonnet-4-6 & Best reasoning/cost balance at simulation scale \\
Structured output & Anthropic tool\_use & Guarantees JSON schema compliance \\
Logging & SQLite (stdlib) & No dependencies, full replay, queryable \\
CLI display & Rich & Turn-by-turn terminal output \\
Dependencies & uv + pyproject.toml & Fast, modern Python dependency management \\
\bottomrule
\caption{Technology stack}
\end{longtable}

% ============================================================
\section{Limitations and Risks}
% ============================================================

\subsection{Named Actor Bias}

WarAgent (2023) demonstrated that naming actors (``China,'' ``USA'') causes LLMs to retrieve training-data priors rather than reason from world state. OSE's Taiwan Strait scenario uses named actors for interpretability, but this introduces training prior contamination. Mitigation: a secondary experimental condition with anonymized actors (``Rising Maritime Power,'' ``Island Democracy'') will test whether actor names affect doctrine fidelity scores.

\subsection{Escalation Bias}

Every prior study finds LLMs escalate systematically. OSE's inaction action category (hold\_position, monitor, delay\_commitment) and explicit cost penalties for escalatory actions partially mitigate this. However, if all 80 runs produce war outcomes regardless of doctrine, the experimental design loses discrimination power. Mitigation: initial parameter tuning on 5 runs before full experimental execution.

\subsection{LLM Variance}

Temperature=0 reduces but does not eliminate output variance. The same prompt can produce different outputs across API calls due to non-determinism in sampling infrastructure. Mitigation: full prompt logging for replay analysis; behavioral consistency scores measure variance explicitly rather than assuming it away.

\subsection{Doctrine Scoring Validity}

Doctrine Fidelity Scoring uses a secondary LLM call to evaluate reasoning traces. This introduces measurement error: the scorer may disagree with ground truth on marginal cases. Mitigation: scorer rubric is explicit and structured; sample of 100 decisions will be manually reviewed against scorer output to validate calibration.

\subsection{External Validity}

OSE's findings apply to LLM behavior in a controlled simulation environment. Generalization to real LLM deployments in actual strategic analysis contexts requires caution. The simulation abstracts away many factors present in real deployments (classified information, human-AI teaming dynamics, institutional context). OSE's findings are indicative, not directly transferable.

% ============================================================
\section{Open Questions}
% ============================================================

\begin{enumerate}[leftmargin=1.5em]
    \item \textbf{Doctrine vs. personality:} Is doctrine-following behavior model-specific (Claude can follow realist doctrine; Gemini cannot) or general? Secondary comparison with a second model is planned but not yet scoped.

    \item \textbf{Doctrine specification granularity:} How detailed does the doctrine specification need to be to produce measurable fidelity? A five-step procedure vs. a one-paragraph description---does granularity matter?

    \item \textbf{Cross-run consistency measurement:} How many runs are needed to establish reliable BCI estimates? Statistical power analysis suggests 20 runs per condition is sufficient for detecting large effects but may underpower detection of moderate effects.

    \item \textbf{Communication phase effects:} Does the free-form communication phase improve or degrade doctrine fidelity? Actors may express doctrine-consistent reasoning in communication but make off-doctrine decisions. Analyzing communication content against action selection may reveal misalignment between stated and actual doctrine adherence.

    \item \textbf{Scorer calibration:} What is the inter-rater reliability between the LLM fidelity scorer and human raters on the same reasoning traces? This validation step is essential before publishing DFS results.
\end{enumerate}

% ============================================================
\section{Conclusion}
% ============================================================

The Omni-Simulation Engine represents a disciplined attempt to answer a question with genuine academic and practical significance: can large language models make strategic decisions consistent with prescribed state decision doctrines, and how trustworthy can they be for such purposes?

OSE's contribution rests on three pillars. First, a research design that has not been executed before---varying qualitative IR decision doctrine as the experimental treatment in LLM conflict simulation. Second, an architectural innovation that has not been implemented before---a formal perception filter that operationalizes Jervis's misperception thesis as a designed, tunable mechanism. Third, a measurement framework that produces specific, actionable answers to the trustworthiness question rather than general observations about LLM behavior.

The findings will be one of three things: doctrines work, doctrines fail under pressure, or doctrines partially work with specific failure modes. Any of these outcomes is a real finding with real implications---for IR theory, for LLM capability research, and for the policy and defense institutions currently trying to determine how much to trust these systems in strategic contexts.

The project starts narrow: one scenario, one LLM, four doctrine conditions, twenty runs each. If the core loop works and the measurement framework produces coherent results, scope expands. The ambition is large. The method must remain disciplined.

% ============================================================
\section*{References}
% ============================================================
\addcontentsline{toc}{section}{References}

\begin{enumerate}[leftmargin=1.5em, label={[\arabic*]}]
    \item Allison, G. (1971). \textit{Essence of Decision: Explaining the Cuban Missile Crisis}. Little, Brown.

    \item Bakhtin, A., et al. (Meta AI). (2022). Human-level play in the game of Diplomacy by combining language models with strategic reasoning. \textit{Science}, 378(6624), 1067--1074.

    \item Fearon, J. D. (1994). Domestic political audiences and the escalation of international disputes. \textit{American Political Science Review}, 88(3), 577--592.

    \item Ge, X., et al. (2023). War and Peace (WarAgent): Large Language Model-based Multi-Agent Simulation of World Wars. \textit{arXiv:2311.17227}.

    \item Jervis, R. (1976). \textit{Perception and Misperception in International Politics}. Princeton University Press.

    \item Keohane, R. O. (1984). \textit{After Hegemony: Cooperation and Discord in the World Political Economy}. Princeton University Press.

    \item Li, M., Li, X., \& Zhou, T. (2026). When AI Navigates the Fog of War. \textit{arXiv:2603.16642}.

    \item Mearsheimer, J. J. (2001). \textit{The Tragedy of Great Power Politics}. Norton.

    \item Moghimi, S., et al. (2024). Escalation Risks from Language Models in Military and Diplomatic Decision-Making. \textit{Proceedings of the ACM Conference on Fairness, Accountability, and Transparency}. arXiv:2401.03408.

    \item Park, J. S., et al. (2023). Generative Agents: Interactive Simulacra of Human Behavior. \textit{Proceedings of UIST 2023}. arXiv:2304.03442.

    \item Payne, K., et al. (2026). AI Arms and Influence: Frontier Models Exhibit Sophisticated Reasoning in Simulated Nuclear Crises. \textit{arXiv:2602.14740}.

    \item Raz, I., et al. (2025). Prospect Theory Fails for LLMs: Revealing Instability of Decision-Making under Epistemic Uncertainty. \textit{arXiv:2508.08992}.

    \item Schelling, T. C. (1960). \textit{The Strategy of Conflict}. Harvard University Press.

    \item Schelling, T. C. (1966). \textit{Arms and Influence}. Yale University Press.

    \item Waltz, K. N. (1979). \textit{Theory of International Politics}. Addison-Wesley.

    \item Xu, Z., et al. (2024). Open-Ended Wargames with Large Language Models. \textit{arXiv:2404.11446}.

    \item Yang, X., et al. (2025). MultiAgentBench: Evaluating the Collaboration and Competition of LLM Agents. \textit{arXiv:2503.01935}.

    \item Zhang, Y., et al. (2026). LLMs as Strategic Actors: Behavioral Alignment, Risk Calibration, and Argumentation Framing in Geopolitical Simulations. \textit{arXiv:2603.02128}.

    \item Zhu, K., et al. (2024). LLM-Deliberation: Evaluating LLMs with Interactive Multi-Agent Negotiation Games. \textit{OpenReview}.
\end{enumerate}

\vspace{2cm}
\hrule
\vspace{0.5cm}
{\small \textit{Omni-Simulation Engine --- OSE Research Document --- Daniel Reshetnikov --- March 2026}}

\end{document}
